\documentclass[10pt,a4paper]{amsart}
\usepackage[margin=19mm]{geometry}
\usepackage{amsmath,amssymb,mathtools}
\usepackage[scr=rsfs,cal=euler]{mathalfa}
\usepackage{xfrac}
\usepackage[unicode,hidelinks,bookmarks=false]{hyperref}
\newcommand{\ie}{i.e.\ }
\newcommand{\cf}{cf.\ }
\newcommand{\resp}{respectively\ }
\newcommand{\Subsection}{\subsection}
\providecommand{\nobreakdash}{\nobreak\hbox{-}}
\setlength{\emergencystretch}{2em}
\multlinegap0pt
\usepackage{bm}
\usepackage{bbm}
\usepackage{dsfont}
\usepackage{xspace}
\usepackage{cancel}
\usepackage{mathtools}
\usepackage{amsthm}
\makeatletter
\@ifundefined{theorem}{\newtheorem{theorem}{Theorem}}{}
\@ifundefined{corollary}{\newtheorem{corollary}[theorem]{Corollary}}{}
\@ifundefined{proposition}{\newtheorem{proposition}[theorem]{Proposition}}{}
\makeatother
\DeclareMathOperator{\GL}{GL}
\DeclareMathOperator{\PGL}{PGL}
\DeclareMathOperator{\PSL}{PSL}
\DeclareMathOperator{\SL}{SL}
\DeclareMathOperator{\nrd}{nrd}
\DeclareMathOperator{\ram}{Ram}
\title[An Arithmetic Invariant of the Jacquet-Langlands corresponde]{An Arithmetic Invariant of the Jacquet-Langlands correspondence}
\author{Jun Yang}
\date{Source: arXiv:2405.18372v5 (CC BY 4.0)}
\begin{document}
\maketitle

\noindent\emph{Source and licence.} An Arithmetic Invariant of the Jacquet-Langlands correspondence. Version arXiv:2405.18372v5 (CC BY 4.0), distributed under the Creative Commons Attribution 4.0 International licence, as recorded in the arXiv licence metadata for this exact version and shown on the abstract page https://arxiv.org/abs/2405.18372v5. Full rights notice: licenses/arxiv-2405.18372-CC-BY-4.0.txt.\par\smallskip

\noindent\textit{Editorial scope.} Counted unit: the Corollary whose source label is \texttt{cPGLhaar=} in the article (source0.tex lines 1290--1293, source label \texttt{cPGLhaar=}), delivered together with the article's own complete original proof of it (source0.tex lines 1294--1338, one \texttt{proof} environment of 45 lines). No mathematics is written, repaired or re-derived: statement and proof are the article's own text line for line. Required context retained uncounted: pStama (lines 1253--1262). Every definition, display and label of those lines is retained unchanged. Omitted from this package and not redistributed: 1--1252, 1263--1289, 1339--1430. Nothing inside a retained excerpt is omitted or altered: every retained body is delivered exactly as the article's lines are written, so no omission receipt of any kind accompanies this package. Citations: the retained excerpts carry no citation command at all, so no bibliography entry of the article is transcribed and no reference is redistributed. Reference adaptation: none was needed and none was made; every reference inside the delivered unit resolves inside it, so no unresolved reference or citation is produced. Printed slips kept literal: no printed word, symbol, hypothesis or number of the counted statement or of its proof was altered. Layout and class: the article's own class is \texttt{article} with packages including geometry, ifsym, lipsum, amssymb, amsmath, latexsym, bm, enumitem, graphicx, amscd, extarrows, amsfonts, indentfirst, bbm; the wrapper is the lane's pinned \texttt{amsart} editorial wrapper at 10pt on A4, which restates the theorem-like environments the retained text needs and adds the article's own macro definitions that the retained text needs (\texttt{\textbackslash GL}, \texttt{\textbackslash PGL}, \texttt{\textbackslash PSL}, \texttt{\textbackslash SL}, \texttt{\textbackslash nrd}, \texttt{\textbackslash ram}), with extra wrapper packages (bm, bbm, dsfont, xspace, cancel, mathtools). The delivered numbering is the unit's own. Assets: no retained span contains a figure environment, a graphics-inclusion command, a PDF-page inclusion, a table or a TikZ picture; no image file is copied into this package, the package declares no asset dependency and no image file is redistributed with it. Licence: version arXiv:2405.18372v5 of this article is distributed under the Creative Commons Attribution 4.0 International licence, as recorded in the arXiv licence metadata for this exact version and shown on the abstract page https://arxiv.org/abs/2405.18372v5. A full rights notice is delivered at licenses/arxiv-2405.18372-CC-BY-4.0.txt. Characters: every retained excerpt is pure ASCII, so the delivered engine, XeLaTeX with its default Latin Modern text font, reports no missing character.
\begin{proposition}\label{pStama}
Let $G_1,G_2$ be simply connected semisimple algebraic groups over a number field $F$ such that they are inner forms of each other. 
Suppose $\mu_{1,v},\mu_{2,v}$ are compatible Haar measures on $G_1(F_v),G_2(F_v)$. 
Let $\mu_{1,S}=\prod_{v\in S}\mu_{1,v}$ and $\mu_{2,S}=\prod_{v\in S}\mu_{2,v}$. 
We have
\begin{center}
    $\mu_{1,S}(G_{1}(F_S)/\Gamma_{1}(S))=\mu_{2,S}( G_{2}(F_S)/\Gamma_{2}(S))$
\end{center}
if $\ram(G_1)\cup \ram(G_2)\subset S$ and both $G_1(F_S),G_2(F_S)$ are non-compact.  
\end{proposition}

\begin{corollary}\label{cPGLhaar=}
If $\ram(\overline{G'})\subset S$, 
$\mu_{S}(\PGL(n,F_S)/ \PGL(n,\mathcal{O}_S))=\mu'_{S}(\overline{G'}(F_S)/ \overline{G'}(\mathcal{O}_S))$
\end{corollary}

\begin{proof}
Let $R$ be a commutaive ring and consider the following short exact sequence
\begin{center}
    $1\xrightarrow[]{} \SL(n,R)\xrightarrow[]{}\GL(n,R)\xrightarrow[]{\eta}  R^{\times}\to 1$,
\end{center}
where $\eta$ is the determinant map. 
Let $\overline{\eta}$ denote the induced map $\PGL(n,R)\to R^{\times}/(R^{\times})^{n}$. 
It is known that $\ker \overline{\eta}=\PSL(n,R)$ and
\begin{center}
    $\PGL(n,R)/\PSL(n,R)\cong R^{\times}/(R^{\times})^{n}$. 
\end{center}
On the other hand, $\SL(n,R)/\mu_n(R)\cong \PSL(n,R)$, where $\mu_{n}(R)$ is the group of $n$-th roots of unity in $\mathbb{R}^{\times}$.

Let $\omega,\omega'$ be the invariant differential forms of top degree on $\overline{G}(F)=\PGL(n,F),\overline{G'}$ which determines the compatible local Haar measures $\mu_v,\mu'_{v}$. 

We first consider the split case, i.e., $\overline{G}=\PGL(n)$. 
Note that both the maps $\SL(n,R)\to \PSL(n,R)$ and $\PGL(n,R)\to \PSL(n,R)$ give coverings of finite index when $R$ is a local field. 
We may pull back $\omega$ to the larger groups. It gives compatible Haar measures on $\SL(n,F_v),\PSL(n,F_v)$, which we will also denote by $\mu_{v}$.
Thus, we obtain
\begin{equation}\label{ePGL-SL}
\begin{aligned}
    \mu_{S}(\PGL(n,F_S)/\PGL(n,\mathcal{O}_S))&=\frac{[\PGL(n,F_S):\PSL(n,F_S)]}{[\PGL(n,\mathcal{O}_S):\PSL(n,\mathcal{O}_S)]}\cdot \mu_{S}(\PSL(n,F_S)/\PSL(n,\mathcal{O}_S))\\
    &=\frac{[F_S^{\times}:(F_S^{\times})^n]}{[\mathcal{O}_S^{\times}:(\mathcal{O}_S^{\times})^n]}\cdot \frac{1}{[\mu_n(F_S):\mu_n(\mathcal{O}_S)]}\cdot \mu_{S}(\SL(n,F_S)/\SL(n,\mathcal{O}_S)). 
\end{aligned}
\end{equation}

For the non-split group $G'$ and $\overline{G'}=G'/Z(G')$, we have the reduced norm $\nrd$ which gives the short exact sequence
\begin{center}
    $1\xrightarrow[]{} G'^{(1)}(R)\xrightarrow[]{}G'(R)\xrightarrow[]{\nrd}  R^{\times}\to 1$,
\end{center}
where $G'^{(1)}=\{g\in G'|\nrd(g)=1\}$. 
Let $\overline{\nrd}$ denote the induced map on $(G'/Z(G'))(R)\to R^{\times}/(R^{\times})^{n}$.
We denote $(G'^{(1)}/Z(G'^{(1)}))$ by $\overline{G'^{(1)}}$. 
We know that $\ker\nrd=\overline{G'^{(1)}}(R)$, which is an inner form of $\PSL(n,R)$ when $R$ is a local field. 


On the other hand, we have $G'^{(1)}(R)/\mu_{n}(R)\cong \overline{G'^{(1)}}(R)$. 
The same arguments as in Equation \eqref{ePGL-SL} give
\begin{equation}\label{ePG'G'1}
\begin{aligned}
    \mu'_{S}(\overline{G'}(n,F_S)/\overline{G'}(n,\mathcal{O}_S))&=\frac{[\overline{G'}(n,F_S):\overline{G'^{(1)}}(n,F_S)]}{[\overline{G'}(n,\mathcal{O}_S):\overline{G'^{(1)}}(n,\mathcal{O}_S)]}\cdot \mu_{S}(\overline{G'^{(1)}}(n,F_S)/\overline{G'^{(1)}}(n,\mathcal{O}_S))\\&=\frac{[F_S^{\times}:(F_S^{\times})^n]}{[\mathcal{O}_S^{\times}:(\mathcal{O}_S^{\times})^n]}\cdot \frac{1}{[\mu_n(F_S):\mu_n(\mathcal{O}_S)]}\cdot \mu'_{S}(G'^{(1)}(n,F_S)/G'^{(1)}(n,\mathcal{O}_S)). 
\end{aligned}
\end{equation}
By Theorem \ref{pStama}, the comparison of Equations \eqref{ePGL-SL} and \eqref{ePG'G'1} gives the desired formula. 
\end{proof}

\end{document}
